%---------- Inleiding ---------------------------------------------------------

\section{Inleiding}%
\label{sec:inleiding}
Cybercriminaliteit die gericht is op het verkrijgen van geld of persoonsgegevens is geen recent verschijnsel. Reeds in de jaren negentig slaagden aanvallers erin om via phishing de authenticatiegegevens van America Online (AOL)-gebruikers te bemachtigen \autocite{Goel2018}. Sindsdien heeft phishing zich verder ontwikkeld en is het uitgegroeid tot een veelvoorkomende vorm van cybercriminaliteit. De blijvende relevantie van deze vorm van criminaliteit blijkt onder meer uit recente cijfers over de omvang van phishing in België.

In 2025 stuurden burgers bijna 10 miljoen verdachte berichten door naar \url{verdacht@safeonweb.be}\footnote{Centrum voor Cybersecurity België, \url{https://ccb.belgium.be/nl}, geraadpleegd op 8 oktober 2026.}. De omvang van phishing is bovendien niet uitsluitend zichtbaar in het aantal gemelde berichten, maar ook in de financiële schade die ermee gepaard gaat. Volgens \textcite{Febelfin2025} maakten phishers in 2024 ongeveer 49 miljoen euro buit.

Hoewel phishing via verschillende kanalen verloopt, blijft misleiding via e-mail volgens \textcite{Veit2025} een van de grootste beveiligingsrisico's voor zowel bedrijven als eindgebruikers. Wanneer zo'n phishingmail de filters van de mailclient passeert, is het de verantwoordelijkheid van de gebruiker om in te schatten of de mail te vertrouwen is of niet. Deze beoordeling wordt bovendien bemoeilijkt door de opkomst van generatieve artificiële intelligentie, die aanvallers in staat stelt om overtuigende en gepersonaliseerde phishingberichten te creëren \autocite{SchmittFlechais2024}.

Hoewel sensibilisering gebruikers kan helpen om phishingmails te herkennen, volstaat ze niet om het risico op succesvolle aanvallen volledig weg te nemen. Onderzoek toont aan dat gebruikers ook na phishing\-training kwetsbaar blijven voor dergelijke aanvallen \autocite{Marshall2024}. 
Daarom is aanvullende ondersteuning wenselijk op het moment dat de gebruiker een e-mail moet beoordelen. 

Deze bachelorproef onderzoekt daarom de volgende centrale onderzoeksvraag: \emph{``Hoe kunnen relevante kenmerken uit e-mails worden geïdentificeerd en gecombineerd om phishingmails automatisch te detecteren en hun risiconiveau te bepalen?''}

Naast de scriptie levert het onderzoek een algoritme op dat een e-mail analyseert, er een risicoscore aan toekent en de onderliggende redenen op een begrijpelijke manier aan de gebruiker toont. Om de centrale onderzoeksvraag te beantwoorden, worden de volgende deelvragen behandeld:

\begin{enumerate}
  \item \emph{Welke kenmerken zijn het meest relevant om phishing te detecteren?}
  \item \emph{Hoe kunnen deze kenmerken automatisch uit een e-mail worden gehaald?}
  \item \emph{Hoe kunnen deze kenmerken worden gecombineerd tot een risicoscore?}
\end{enumerate}

Het onderzoek wordt als geslaagd beschouwd wanneer het algoritme, geëvalueerd op een vooraf afgebakende dataset van phishingmails en legitieme e-mails, phishing betrouwbaar van legitieme e-mail onderscheidt, gemeten aan de hand van evaluatiecriteria die in de literatuurstudie worden onderbouwd.

%---------- Stand van zaken ---------------------------------------------------

\section{Literatuurstudie}%
\label{sec:literatuurstudie}

Hier beschrijf je de \emph{state-of-the-art} rondom je gekozen onderzoeksdomein, d.w.z.\ een inleidende, doorlopende tekst over het onderzoeksdomein van je bachelorproef. Je steunt daarbij heel sterk op de professionele \emph{vakliteratuur}, en niet zozeer op populariserende teksten voor een breed publiek. Wat is de huidige stand van zaken in dit domein, en wat zijn nog eventuele open vragen (die misschien de aanleiding waren tot je onderzoeksvraag!)?

Je mag de titel van deze sectie ook aanpassen (literatuurstudie, stand van zaken, enz.). Zijn er al gelijkaardige onderzoeken gevoerd? Wat concluderen ze? Wat is het verschil met jouw onderzoek?

Verwijs bij elke introductie van een term of bewering over het domein naar de vakliteratuur, bijvoorbeeld~\autocite{Hykes2013}! Denk zeker goed na welke werken je refereert en waarom.

Draag zorg voor correcte literatuurverwijzingen! Een bronvermelding hoort thuis \emph{binnen} de zin waar je je op die bron baseert, dus niet er buiten! Maak meteen een verwijzing als je gebruik maakt van een bron. Doe dit dus \emph{niet} aan het einde van een lange paragraaf. Baseer nooit teveel aansluitende tekst op eenzelfde bron.

Als je informatie over bronnen verzamelt in JabRef, zorg er dan voor dat alle nodige info aanwezig is om de bron terug te vinden (zoals uitvoerig besproken in de lessen Research Methods).

% Voor literatuurverwijzingen zijn er twee belangrijke commando's:
% \autocite{KEY} => (Auteur, jaartal) Gebruik dit als de naam van de auteur
%   geen onderdeel is van de zin.
% \textcite{KEY} => Auteur (jaartal)  Gebruik dit als de auteursnaam wel een
%   functie heeft in de zin (bv. ``Uit onderzoek door Doll & Hill (1954) bleek
%   ...'')

Je mag deze sectie nog verder onderverdelen in subsecties als dit de structuur van de tekst kan verduidelijken.

%---------- Methodologie ------------------------------------------------------
\section{Methodologie}%
\label{sec:methodologie}

Hier beschrijf je hoe je van plan bent het onderzoek te voeren. Welke onderzoekstechniek ga je toepassen om elk van je onderzoeksvragen te beantwoorden? Gebruik je hiervoor literatuurstudie, interviews met belanghebbenden (bv.~voor requirements-analyse), experimenten, simulaties, vergelijkende studie, risico-analyse, PoC, \ldots?

Valt je onderwerp onder één van de typische soorten bachelorproeven die besproken zijn in de lessen Research Methods (bv.\ vergelijkende studie of risico-analyse)? Zorg er dan ook voor dat we duidelijk de verschillende stappen terug vinden die we verwachten in dit soort onderzoek!

Vermijd onderzoekstechnieken die geen objectieve, meetbare resultaten kunnen opleveren. Enquêtes, bijvoorbeeld, zijn voor een bachelorproef informatica meestal \textbf{niet geschikt}. De antwoorden zijn eerder meningen dan feiten en in de praktijk blijkt het ook bijzonder moeilijk om voldoende respondenten te vinden. Studenten die een enquête willen voeren, hebben meestal ook geen goede definitie van de populatie, waardoor ook niet kan aangetoond worden dat eventuele resultaten representatief zijn.

Uit dit onderdeel moet duidelijk naar voor komen dat je bachelorproef ook technisch voldoen\-de diepgang zal bevatten. Het zou niet kloppen als een bachelorproef informatica ook door bv.\ een student marketing zou kunnen uitgevoerd worden.

Je beschrijft ook al welke tools (hardware, software, diensten, \ldots) je denkt hiervoor te gebruiken of te ontwikkelen.

Probeer ook een tijdschatting te maken. Hoe lang zal je met elke fase van je onderzoek bezig zijn en wat zijn de concrete \emph{deliverables} in elke fase?

%---------- Verwachte resultaten ----------------------------------------------
\section{Verwacht resultaat, conclusie}%
\label{sec:verwachte_resultaten}

Hier beschrijf je welke resultaten je verwacht. Als je metingen en simulaties uitvoert, kan je hier al mock-ups maken van de grafieken samen met de verwachte conclusies. Benoem zeker al je assen en de onderdelen van de grafiek die je gaat gebruiken. Dit zorgt ervoor dat je concreet weet welk soort data je moet verzamelen en hoe je die moet meten.

Wat heeft de doelgroep van je onderzoek aan het resultaat? Op welke manier zorgt jouw bachelorproef voor een meerwaarde?

Hier beschrijf je wat je verwacht uit je onderzoek, met de motivatie waarom. Het is \textbf{niet} erg indien uit je onderzoek andere resultaten en conclusies vloeien dan dat je hier beschrijft: het is dan juist interessant om te onderzoeken waarom jouw hypothesen niet overeenkomen met de resultaten.

